\chapter{Checks}
\label{ch:checks}

\cmd{bag check} reads the workbook and runs every check before anything is written.
\cmd{bag ir}, \cmd{bag generate} and \cmd{bag flow} run the same checks first, on a
workbook and on an IR alike (chapter~\ref{ch:ir}).

\section{How a finding reads}

Each problem the checks find is a \term{finding}. It is one line:

\begin{center}
\code{<severity>: <Sheet> row <n>: <what is wrong> (<RULE-ID>)}
\end{center}

\begin{description}
\item[severity] \code{error} or \code{warning}. An error stops BAG: \cmd{bag check}
  exits with status 1, and nothing is generated. A warning is printed and BAG goes on.
\item[sheet and row] Where to look in the workbook. The row is the Excel row number.
  A finding about something that is missing has no row, for example
  \code{Address_Decode: master port dma0 on main_xbar has no row}.
\item[what is wrong] The value, and what to do about it where that is not obvious.
\item[rule ID] A stable name for the check, such as \ruleid{TOPO-INDEX-OUT-OF-RANGE}.
  Its first word names the area: \code{TOPO-} the Topology sheet, \code{MM-} the memory
  map, \code{DEC-} address decode, and so on. Use the rule ID to find the check in the
  table below.
\end{description}

After the findings, \cmd{bag check} prints one summary line, for example
\code{spec.xlsx: 0 error(s), 2 warning(s)}.

Findings come in the order of the sheets and rows, so fix the first ones first: one
mistake can cause later findings. For example, a misspelt fabric name on the Topology
sheet also leaves ports unconnected.

\section{Two kinds of check}

Most checks look at one sheet, or compare sheets: that a name is defined, that a value
is legal, that two regions do not overlap. A few checks need the resolved design, for
example to know which adapters BAG inserts. They run only when the others find no error.
The two \ruleid{AXI-WRAP-REFUSED} warnings of the sample (chapter~\ref{ch:first}) are of
this kind: they follow
the path from a master through its adapters.

A few problems are found only by \cmd{bag generate}. It then stops with \code{error:}
and a message that names the fabric, instance or name, and writes no RTL. These are:

\begin{itemize}
\item two names that become the same name in the generated files: two entries of the
  parameter package (\code{package names collide}), two C macros, a generated net and
  a wrapper port, or a sub-wrapper with the name of the top wrapper. Rename one of them.
\item two user-supplied fabrics that share an External module but would need different
  ports, for example through different ID widths.
\item a fabric whose masters see different address translations: a pass-through N
  bridge on the same fabric as a direct master, or as a bridge with another window. One
  decoder serves all the masters of a fabric, so give them separate fabrics.
\end{itemize}

\section{Every check}

The table lists every rule, in alphabetical order, with its severity and the message it
gives on a broken copy of the sample workbook. The table is generated from the checks
themselves, so it is complete. Your message names your own sheet rows, ports and values.

\begingroup
\footnotesize
\begin{longtable}{@{}>{\raggedright}p{0.28\linewidth}>{\raggedright}p{0.08\linewidth}p{0.58\linewidth}@{}}
\toprule
Rule ID & Severity & Example message \tabularnewline
\midrule
\endfirsthead
\toprule
Rule ID & Severity & Example message \tabularnewline
\midrule
\endhead
% \input cannot run inside a table (it leaves \bottomrule misplaced); the primitive can.
\csname @@input\endcsname checks.tex
\bottomrule
\end{longtable}
\endgroup
